% Generated draft: requires mathematical and source review.
\documentclass[11pt]{article}
\usepackage[margin=25mm]{geometry}
\usepackage{fontspec}
\setmainfont{Latin Modern Roman}
\usepackage{amsmath,amssymb,mathtools,unicode-math}
\setmathfont{Latin Modern Math}
\usepackage{enumitem,needspace,xurl,hyperref}
\setlength{\parindent}{0pt}
\setlength{\parskip}{5pt}
\setlength{\emergencystretch}{3em}
\newcommand{\N}{\mathbb{N}}
\newcommand{\Z}{\mathbb{Z}}
\newcommand{\Q}{\mathbb{Q}}
\newcommand{\R}{\mathbb{R}}
\newcommand{\C}{\mathbb{C}}
\newcommand{\F}{\mathbb{F}}
\newcommand{\A}{\mathbb{A}}
\newcommand{\PP}{\mathbb{P}}
\newcommand{\E}{\mathbb{E}}
\newcommand{\eps}{\varepsilon}
\newcommand{\dd}{\,\mathrm{d}}
\begin{document}
\section*{UnsolvedMath ID: 700926}
\textbf{Source problem code:} catalogue entry
\subsection{Definitions and mathematical statement}
A group \(G\) is \emph{periodic} if every element has finite order, \emph{locally soluble} if every finitely generated subgroup is soluble, and \emph{locally nilpotent} if every finitely generated subgroup is nilpotent. For \(g\in G\), its centralizer is
\[
C_G(g)=\{x\in G:xg=gx\}.
\]
A group \(H\) is \emph{Chernikov} if it has a normal subgroup \(A\) of finite index such that
\[
A\cong\prod_{j=1}^{r}C_{p_j^\infty},
\]
where \(r\geq 0\), each \(p_j\) is a prime, and \(C_{p^\infty}\) is the Prüfer \(p\)-group, the union of an ascending chain of cyclic groups of orders \(p,p^2,\ldots\). The empty product is trivial, so finite groups are Chernikov.

For each of the following hypotheses separately, the question asks whether every locally soluble periodic group \(G\) satisfying that hypothesis admits a finite normal series
\[
1=G_0\trianglelefteq G_1\trianglelefteq\cdots\trianglelefteq G_n=G
\]
for some integer \(n\geq 0\), with \(G_{i+1}/G_i\) locally nilpotent for every \(0\leq i<n\):
\begin{enumerate}
\item[(a)] There exists \(g\in G\) such that \(C_G(g)\) is finite.
\item[(b)] There exists \(g\in G\) such that \(C_G(g)\) is Chernikov.
\end{enumerate}
There is no prescribed uniform bound on the length \(n\). The displayed chain specifies successive normality. The supplied wording does not explicitly settle whether every \(G_i\) must additionally be normal in \(G\); that convention remains a formulation issue for review.

The supplied registry labels the problem partially solved and reports an affirmative solution of part (a), attributed to E. Khukhro. This is a supplied source assessment, not an independently verified mathematical assertion here. No resolution of part (b) is supplied.
\subsection{Short English statement}
For a locally soluble group in which every element has finite order, does the existence of an element with finite centralizer, or more generally Chernikov centralizer, force a finite normal series with locally nilpotent factors? The supplied record reports an affirmative solution of part (a); it supplies no resolution of part (b). The intended normality convention for the series needs review.
\subsection{Sources}
\begin{itemize}
\item Supplied local UnsolvedMath registry record: permanent ID \(700926\), exact source code \texttt{catalogue entry}, title \emph{Kourovka Notebook Problem 13.8}. The registry lists V. V. Belyaev as proposer and gives the year 1995; the statement itself carries the attribution B. Hartley. Both attributions are retained without inferring their relationship.
\item \emph{The Kourovka Notebook (Unsolved Problems in Group Theory)}, problem 13.8, 13th issue (1995), as represented in the supplied record by the 21st issue edition (2026), \texttt{arXiv:1401.0300v48}. The supplied collection description names E. I. Khukhro and V. D. Mazurov as editors.
\item Supplied discussion reference for the affirmative resolution of part (a): E. Khukhro, \emph{Bull. London Math. Soc.}, volume 58, no. 1 (2026), Article ID e70206. No article title or direct article link is supplied.
\item Supplied status review dated \texttt{2026-10-16}: canonical status \texttt{partially\_solved}, evidence category \texttt{maintained\_source}. Its note reports that part (a) is marked solved affirmatively in the October 2026 21st issue, \texttt{arXiv:1401.0300v48}. This assessment is supplied dated curation and has not been independently verified.
\end{itemize}
\end{document}
